\documentclass[11pt]{article}
\usepackage[utf8]{inputenc}
\usepackage{amsmath}
\usepackage{amssymb}
\usepackage{geometry}
\usepackage{enumitem}
\usepackage{titlesec}
\usepackage{xcolor}
\usepackage{hyperref}
\usepackage{multicol}
\usepackage{array}

\geometry{a4paper, margin=0.75in}
\hypersetup{colorlinks=true, linkcolor=blue!70!black, urlcolor=blue}

% Colored boxes for formula cheat sheets
\usepackage{tcolorbox}
\tcbuselibrary{skins}
\newtcolorbox{formulabox}[1][]{
  colback=blue!5!white,
  colframe=blue!60!black,
  fonttitle=\bfseries,
  title=#1,
  arc=3pt,
  boxrule=0.8pt
}

\title{\textbf{EECS 498/598-016 Final Exam Study Guide}\\
\large{Winter 2026 - Conceptual Reference with Key Formulas}}
\author{Based on Lecture Slides + Exam Prep}
\date{\today}

\begin{document}
\maketitle
\vspace{-1em}

\begin{abstract}
\noindent Exam format: 70 minutes, closed-note, 6 questions. Focus on understanding core concepts, not proofs, programming, or complex equations. \textbf{Key formulas are marked with \colorbox{blue!10}{\textcolor{blue!60!black}{Formula Box}}.}
\end{abstract}

\tableofcontents
\newpage

\section{Rasterization}
\subsection{Transformations \& Projections}
\begin{itemize}
    \item \textbf{Homogeneous Coordinates}: Enable translation to be represented as matrix multiplication (same as rotation/scaling). 2D points become $(x, y, 1)$, allowing any affine transform to use a $3\times3$ matrix. 3D points become $(x, y, z, 1)$ with $4\times4$ matrices.
    \item \textbf{Composition}: Chain transforms via matrix multiplication. Order matters (non-commutative: $M_1 M_2 \neq M_2 M_1$).
    \item \textbf{Viewing Transform}: Camera matrix positions/orients the world relative to viewer.
    \item \textbf{Projection}: 
    \begin{itemize}
        \item \textit{Orthographic}: Parallel projection, preserves relative sizes.
        \item \textit{Perspective}: Converging lines (foreshortening), mimics human eye. Uses frustum and division by $z$ (or $w$).
    \end{itemize}
\end{itemize}

\begin{formulabox}{2D Transformation Matrices}
Scaling (factors $s_x, s_y$):
$$\begin{bmatrix} x' \\ y' \end{bmatrix} = \begin{bmatrix} s_x & 0 \\ 0 & s_y \end{bmatrix} \begin{bmatrix} x \\ y \end{bmatrix}$$

Reflection about $y$-axis:
$$\begin{bmatrix} x' \\ y' \end{bmatrix} = \begin{bmatrix} -1 & 0 \\ 0 & 1 \end{bmatrix} \begin{bmatrix} x \\ y \end{bmatrix}$$

Shear (horizontal shift $a$ at $y=1$):
$$\begin{bmatrix} x' \\ y' \end{bmatrix} = \begin{bmatrix} 1 & a \\ 0 & 1 \end{bmatrix} \begin{bmatrix} x \\ y \end{bmatrix}$$

Rotation (CCW by $\theta$ about origin):
$$\begin{bmatrix} x' \\ y' \end{bmatrix} = \begin{bmatrix} \cos\theta & -\sin\theta \\ \sin\theta & \cos\theta \end{bmatrix} \begin{bmatrix} x \\ y \end{bmatrix}$$
\end{formulabox}

\begin{formulabox}{Homogeneous Coordinates — Full $3\times3$ Matrices}
Translation by $(t_x, t_y)$:
$$\begin{bmatrix} x' \\ y' \\ 1 \end{bmatrix} = \begin{bmatrix} 1 & 0 & t_x \\ 0 & 1 & t_y \\ 0 & 0 & 1 \end{bmatrix} \begin{bmatrix} x \\ y \\ 1 \end{bmatrix}$$

Rotation (CCW by $\theta$):
$$\begin{bmatrix} x' \\ y' \\ 1 \end{bmatrix} = \begin{bmatrix} \cos\theta & -\sin\theta & 0 \\ \sin\theta & \cos\theta & 0 \\ 0 & 0 & 1 \end{bmatrix} \begin{bmatrix} x \\ y \\ 1 \end{bmatrix}$$
\end{formulabox}

\begin{formulabox}{Composition of Transforms}
To apply transform $A$ then $B$:
$$M = B \cdot A$$
Applied to point $\mathbf{p}$:
$$\mathbf{p}' = B \cdot A \cdot \mathbf{p}$$
\emph{Note: rightmost transform is applied first.}
\end{formulabox}

\subsection{Barycentric Coordinates}
\begin{itemize}
    \item Represent any point inside a triangle as a weighted sum of its vertices: $P = \alpha A + \beta B + \gamma C$, where $\alpha+\beta+\gamma=1$.
    \item Point is inside triangle iff $\alpha, \beta, \gamma \in [0, 1]$.
    \item Used for \textbf{interpolation}: smoothly blend vertex attributes (color, normals, texture coordinates) across the triangle surface.
\end{itemize}

\begin{formulabox}{Barycentric Interpolation}
Given vertex attributes $f_A, f_B, f_C$:
$$f_P = \alpha \cdot f_A + \beta \cdot f_B + \gamma \cdot f_C$$
where $\alpha + \beta + \gamma = 1$.\\
Point $P$ is inside $\triangle ABC$ iff $\alpha,\beta,\gamma \in [0, 1]$.
\end{formulabox}

\subsection{Blinn-Phong Shading}
Empirical local illumination model with three components:
\begin{itemize}
    \item \textbf{Diffuse}: Surface scatters light equally in all directions. Depends on angle between light direction $\vec{L}$ and surface normal $\vec{N}$.
    \item \textbf{Specular}: Shiny highlights. \textit{Blinn} variant uses the \textbf{halfway vector} $\vec{H}$.
    \item \textbf{Ambient}: Constant baseline lighting.
\end{itemize}

\begin{formulabox}{Blinn-Phong Model}
Total color at point:
$$I = k_a I_a + k_d I_l \max(0, \vec{N} \cdot \vec{L}) + k_s I_l \max(0, \vec{N} \cdot \vec{H})^n$$
where:
\begin{itemize}
    \item $\vec{H} = \frac{\vec{L} + \vec{V}}{\|\vec{L} + \vec{V}\|}$ (halfway vector)
    \item $\vec{L}$ = light direction, $\vec{V}$ = view direction, $\vec{N}$ = surface normal
    \item $k_a, k_d, k_s$ = ambient/diffuse/specular coefficients
    \item $n$ = shininess exponent (larger = tighter highlight)
\end{itemize}
\end{formulabox}

\subsection{Anti-Aliasing}
\begin{itemize}
    \item \textbf{Problem}: Jaggies/staircase artifacts from undersampling high-frequency edges.
    \item \textbf{SSAA (Supersampling)}: Render at higher resolution, then downsample. Computationally expensive.
    \item \textbf{MSAA (Multisampling)}: Sample multiple points per pixel but only compute shading once per polygon. More efficient.
    \item \textbf{Concept}: Low-pass filter high frequencies before sampling (Nyquist-Shannon).
\end{itemize}

\subsection{Depth Buffer (Z-Buffer)}
\begin{itemize}
    \item For each pixel, store the depth ($z$) of the closest fragment seen so far.
    \item When rendering a new fragment, compare its $z$ with the buffer. Only update if new $z$ is closer.
    \item Solves visibility problem efficiently, order-independent.
\end{itemize}

\subsection{Texturing}
\begin{itemize}
    \item Map 2D image (texture) onto 3D surface using UV coordinates.
    \item \textbf{Mipmapping}: Precompute downsampled texture versions. Sample from appropriate level based on distance/footprint to avoid aliasing.
    \item \textbf{Filtering}: Bilinear (interpolates 4 nearest pixels), trilinear (adds mip-level interpolation), anisotropic (handles oblique viewing angles).
\end{itemize}

\section{Ray Tracing}
\subsection{Ray-Surface Intersection}
\begin{itemize}
    \item Ray equation: $\mathbf{r}(t) = \mathbf{o} + t \cdot \mathbf{d}$, where $\mathbf{o}$ = origin, $\mathbf{d}$ = direction, $t > 0$.
    \item For spheres: solve $|\mathbf{r}(t) - \mathbf{c}|^2 = R^2$ for $t$ (quadratic equation).
\end{itemize}

\begin{formulabox}{Ray Equation}
$$\mathbf{r}(t) = \mathbf{o} + t \cdot \mathbf{d}, \quad t \geq 0$$
Find smallest $t > 0$ where ray intersects an object.
\end{formulabox}

\subsection{The Rendering Equation}
\begin{itemize}
    \item Describes light leaving a point in direction $\omega_o$ as sum of emitted light + reflected light from all incoming directions.
    \item \textbf{Integral}: Requires evaluating contributions from infinitely many directions $\Rightarrow$ solved via Monte Carlo sampling in practice.
    \item Recursive in nature: $L_i$ itself depends on light leaving other surfaces.
\end{itemize}

\begin{formulabox}{Rendering Equation}
$$L_o(\mathbf{x}, \omega_o) = L_e(\mathbf{x}, \omega_o) + \int_{\Omega} f_r(\mathbf{x}, \omega_i, \omega_o) \, L_i(\mathbf{x}, \omega_i) \, (\mathbf{n} \cdot \omega_i) \, d\omega_i$$
\begin{itemize}
    \item $L_o$ = outgoing radiance, $L_e$ = emitted radiance
    \item $f_r$ = BRDF, $L_i$ = incoming radiance
    \item $\mathbf{n} \cdot \omega_i$ = cosine weighting (Lambert's law)
\end{itemize}
\end{formulabox}

\subsection{Recursive Ray Tracing}
\begin{itemize}
    \item Primary rays from camera $\rightarrow$ first hit point.
    \item At each hit, spawn \textbf{reflection rays}, \textbf{refraction rays}, and \textbf{shadow rays} toward lights.
    \item Combine contributions weighted by material properties (reflectivity, transparency).
    \item Stop recursion at max depth or when contribution falls below threshold.
    \item Captures global illumination effects: soft shadows, caustics, color bleeding, interreflections.
\end{itemize}

\subsection{Acceleration Structures}
Naive ray-triangle intersection is $O(N)$. Acceleration reduces to $O(\log N)$:
\begin{itemize}
    \item \textbf{BVH (Bounding Volume Hierarchy)}: Tree structure where each node bounds a subset of geometry with a bounding box. Traverse tree; skip volumes that ray doesn't intersect.
    \item \textbf{Spatial Partitioning}: Divide space into uniform grid or octree. Only test objects in cells the ray passes through.
    \item \textbf{Trade-off}: BVH adapts to geometry density (tight bounds), grids are simpler but may have empty/overfilled cells.
\end{itemize}

\subsection{Light Units \& BRDF}
\begin{itemize}
    \item \textbf{Radiance ($L$)}: Power per unit area per unit solid angle. What cameras/eyes measure. Constant along a ray in vacuum.
    \item \textbf{Irradiance ($E$)}: Power per unit area arriving at a surface. Integral of incoming radiance.
    \item \textbf{BRDF}: $f_r(\omega_i, \omega_o)$ describes how much light from direction $\omega_i$ reflects toward $\omega_o$.
\end{itemize}

\begin{formulabox}{BRDF Properties}
\textbf{Helmholtz reciprocity:}
$$f_r(\omega_i, \omega_o) = f_r(\omega_o, \omega_i)$$
\textbf{Energy conservation:}
$$\int_{\Omega} f_r(\omega_i, \omega_o)(\mathbf{n} \cdot \omega_o)\, d\omega_o \leq 1$$
\end{formulabox}

\section{Geometry \& Representations}
\subsection{3D Representations Comparison}
\begin{itemize}
    \item \textbf{Depth Maps}: 2D array storing distance from camera per pixel. Easy to capture (RGB-D cameras) but only captures visible surfaces from one viewpoint.
    \item \textbf{Point Clouds}: Unordered set of $(x,y,z)$ points. Flexible, memory-efficient for sparse data. No surface connectivity; hard to render smoothly.
    \item \textbf{Meshes}: Vertices + edges + faces (usually triangles). Standard for rendering/modeling. Handles topology well. Fixed resolution; hard to represent fine details without many triangles.
    \item \textbf{Parametric Surfaces}: Defined by mathematical functions (e.g., Bézier, NURBS). Smooth, compact, resolution-independent. Limited to simple topologies.
    \item \textbf{Voxels}: 3D grid of volume elements. Simple, supports boolean operations. Memory intensive ($O(n^3)$), blocky at low resolution.
    \item \textbf{Implicit Surfaces}: Defined by $f(x,y,z)=0$ (e.g., SDFs). Smooth, topology-agnostic, easy to merge. Hard to edit directly; requires marching cubes to extract mesh.
    \item \textbf{Density Fields}: Assign probability/density $\sigma(x,y,z)$ at each point. Foundation for NeRF/volume rendering. Infinite resolution but requires sampling to render.
\end{itemize}

\subsection{Bézier Curves \& Surfaces}
\begin{itemize}
    \item Defined by \textbf{control points}. Curve stays within convex hull of control points.
    \item \textbf{Properties}: Smooth, interpolates endpoints, tangent to first/last control segments.
    \item \textbf{Surfaces}: Tensor product of two Bézier curves. Used in modeling (e.g., car bodies, character design).
\end{itemize}

\begin{formulabox}{Cubic Bézier Curve}
$$P(t) = (1-t)^3 P_0 + 3(1-t)^2 t P_1 + 3(1-t) t^2 P_2 + t^3 P_3, \quad t \in [0, 1]$$
\begin{itemize}
    \item $P_0, P_3$ = endpoints (curve passes through them)
    \item $P_1, P_2$ = control points (shape handles, curve is tangent to $P_0 \rightarrow P_1$ and $P_2 \rightarrow P_3$)
\end{itemize}
\end{formulabox}

\section{Reconstruction}
\begin{formulabox}{Camera Projection}
$$\mathbf{x} = K [R \mid t] \mathbf{X}$$
\begin{itemize}
    \item $K$ = intrinsic matrix $\begin{bmatrix} f_x & s & c_x \\ 0 & f_y & c_y \\ 0 & 0 & 1 \end{bmatrix}$
    \item $[R \mid t]$ = extrinsic matrix (rotation + translation)
    \item $\mathbf{X}$ = 3D world point (homogeneous), $\mathbf{x}$ = 2D image pixel
\end{itemize}
\end{formulabox}

\subsection{Sparse vs. Dense Reconstruction}
\begin{itemize}
    \item \textbf{Sparse}: Recover 3D positions of a \textbf{small set} of distinctive feature points. Fast, used for camera pose estimation and SfM (Structure from Motion).
    \item \textbf{Dense}: Recover 3D geometry for \textbf{every pixel/voxel}. Produces detailed surfaces/meshes. Computationally heavy. Used in MVS, NeRF.
\end{itemize}

\subsection{Camera Calibration}
\begin{itemize}
    \item \textbf{Goal}: Determine intrinsic parameters (focal length, principal point, skew) and extrinsic parameters (rotation, translation relative to world).
    \item Essential for metric reconstruction (true scale/shape) vs. projective reconstruction.
\end{itemize}

\subsection{Epipolar Geometry}
\begin{itemize}
    \item Geometry between two views of the same scene.
    \item \textbf{Epipolar constraint}: A point in one image constrains its match in the other image to lie on a line (the \textbf{epipolar line}).
    \item Reduces 2D search to 1D search for correspondence.
    \item Encoded in the \textbf{Essential Matrix} (calibrated) or \textbf{Fundamental Matrix} (uncalibrated).
\end{itemize}

\begin{formulabox}{Epipolar Constraint}
$$\mathbf{x}'^T F \mathbf{x} = 0$$
\begin{itemize}
    \item $F$ = Fundamental Matrix (uncalibrated), or $E$ = Essential Matrix (calibrated)
    \item $\mathbf{x}, \mathbf{x}'$ = corresponding points in two images
    \item Epipolar line in image 2: $\mathbf{l}' = F \mathbf{x}$
\end{itemize}
\end{formulabox}

\subsection{Eight-Point Algorithm}
\begin{itemize}
    \item Computes the Fundamental/Essential matrix from $\geq 8$ corresponding point pairs.
    \item Linear method: solve homogeneous system $A\mathbf{f} = 0$ via SVD.
    \item \textbf{Key assumption}: Point correspondences are accurate. Sensitive to noise/outliers (often combined with RANSAC).
\end{itemize}

\subsection{Feature Matching \& Multi-View Stereo (MVS)}
\begin{itemize}
    \item \textbf{Feature Matching}: Detect keypoints (SIFT, ORB) across images, match descriptors. Used to establish correspondences for SfM/MVS.
    \item \textbf{MVS}: Given calibrated camera poses, fuse information from multiple views to estimate dense depth/surface.
    \item \textbf{Limits of MVS}: Struggles with textureless surfaces, specular/reflective materials, occlusions, and moving objects. Computationally expensive.
\end{itemize}

\section{Neural Fields}
\subsection{Core Concept \& Benefits}
\begin{itemize}
    \item \textbf{Definition}: Neural network parameterizes a continuous field (e.g., color, density) over spatial coordinates: $F_\theta(\mathbf{x}) \rightarrow \mathbb{R}$.
    \item \textbf{Benefits}: (1) \textit{Resolution-free}: query at arbitrary precision. (2) \textit{Compact}: millions of points stored in MBs of weights. (3) \textit{Differentiable}: enable end-to-end gradient-based optimization.
    \item \textbf{Drawback}: Slow training/inference due to per-point MLP evaluations.
\end{itemize}

\subsection{SDFs \& Marching Cubes}
\begin{itemize}
    \item \textbf{SDF (Signed Distance Function)}: $f(\mathbf{x})$ returns shortest distance to surface (positive outside, negative inside, zero on boundary).
    \item \textbf{Marching Cubes}: Algorithm to extract triangle mesh from volumetric data (like SDF). Divides space into cubes, checks sign at vertices, interpolates edge intersections, generates triangles.
\end{itemize}

\begin{formulabox}{Signed Distance Function}
$$f(\mathbf{x}) \begin{cases}
> 0, & \text{outside surface} \\
= 0, & \text{on surface} \\
< 0, & \text{inside surface}
\end{cases}$$
Surface normal: $\mathbf{n} = \nabla f(\mathbf{x})$ (gradient of SDF)
\end{formulabox}

\subsection{NeRF \& Volume Rendering}
\begin{itemize}
    \item \textbf{NeRF}: Maps $(x,y,z,\theta,\phi) \rightarrow (RGB, \sigma)$ via MLP.
    \item \textbf{Density $\sigma$ over opacity $\alpha$}: $\sigma$ is a continuous property independent of sampling step $\delta_i$. $\alpha_i = 1 - \exp(-\sigma_i \delta_i)$. Allows consistent rendering across different sampling rates.
    \item \textbf{Importance Sampling}: Coarse network predicts density distribution $\rightarrow$ fine network samples more densely near surfaces (high density regions). Improves detail without extra compute.
\end{itemize}

\begin{formulabox}{NeRF Volume Rendering}
Composited color along ray:
$$C = \sum_{i=1}^{N} T_i \alpha_i \mathbf{c}_i, \quad \text{where } T_i = \prod_{j=1}^{i-1} (1 - \alpha_j)$$
Opacity from density:
$$\alpha_i = 1 - \exp(-\sigma_i \delta_i)$$
\begin{itemize}
    \item $\mathbf{c}_i$ = color at sample $i$, $\sigma_i$ = density
    \item $\delta_i$ = distance between samples, $T_i$ = transmittance (accumulated visibility)
\end{itemize}
\end{formulabox}

\subsection{Positional Encoding}
\begin{itemize}
    \item \textbf{Problem}: MLPs struggle to learn high-frequency details (spectral bias).
    \item \textbf{Solution}: Map input coordinates $\mathbf{x}$ to higher-dimensional space using sinusoidal functions.
    \item Enables network to represent fine spatial variations.
\end{itemize}

\begin{formulabox}{Positional Encoding}
$$\gamma(\mathbf{x}) = \big[\sin(2^0 \pi \mathbf{x}), \cos(2^0 \pi \mathbf{x}), \; \sin(2^1 \pi \mathbf{x}), \cos(2^1 \pi \mathbf{x}), \dots, \sin(2^{L-1} \pi \mathbf{x}), \cos(2^{L-1} \pi \mathbf{x})\big]$$
Each coordinate is expanded to $2L$ dimensions via sine/cosine at exponentially increasing frequencies.
\end{formulabox}

\subsection{Gaussian Splatting}
\begin{itemize}
    \item Represents scene as a set of 3D Gaussians (position, covariance, opacity, color).
    \item \textbf{Rendering}: Project Gaussians to 2D, sort by depth, alpha-blend. Extremely fast ($>100$ FPS) because it's rasterization-based, not MLP-based.
    \item Trained via differentiable rendering, similar to NeRF but much faster inference.
\end{itemize}

\subsection{Dynamic Scenes \& Flow Fields}
\begin{itemize}
    \item \textbf{Challenge}: NeRF assumes static scene. Real world has motion.
    \item \textbf{Approach}: Add time $t$ as input to network: $F_\theta(\mathbf{x}, t)$. Or model motion with deformation fields / flow vectors that warp a canonical space over time.
    \item \textbf{Flow Fields}: Predict per-point velocity/displacement to handle non-rigid motion, enabling 4D reconstruction.
\end{itemize}

\section{Generative Models}
\subsection{Discriminative vs. Generative}
\begin{itemize}
    \item \textbf{Discriminative}: Learns boundary between classes. Models $P(y|x)$ (given input, predict label). E.g., classifiers, segmentation networks.
    \item \textbf{Generative}: Learns data distribution to create new samples. Models $P(x)$ or $P(x|y)$. E.g., VAEs, GANs, Diffusion models.
\end{itemize}

\subsection{VAEs vs. GANs}
\begin{itemize}
    \item \textbf{VAE (Variational Autoencoder)}: Encoder maps input to latent distribution $q(z|x)$, decoder samples $z$ to reconstruct.
    \begin{itemize}
        \item \textit{Pros}: Stable training, explicit latent structure, interpretable.
        \item \textit{Cons}: Blurry outputs (MSE loss), limited diversity.
    \end{itemize}
    \item \textbf{GAN (Generative Adversarial Network)}: Generator $G$ creates fakes, Discriminator $D$ distinguishes real vs fake. Minimax game.
    \begin{itemize}
        \item \textit{Pros}: Sharp, high-quality images.
        \item \textit{Cons}: Unstable training, mode collapse, hard to evaluate.
    \end{itemize}
\end{itemize}

\begin{formulabox}{VAE Loss (ELBO) \& GAN Objective}
\textbf{VAE Objective (ELBO):}
$$\mathcal{L} = \mathbb{E}_{q(z|x)}[\log p(x|z)] - \text{KL}(q(z|x) \,\|\, p(z))$$
(reconstruction loss + KL regularization toward $\mathcal{N}(0, 1)$)

\textbf{GAN Minimax:}
$$\min_G \max_D V(D, G) = \mathbb{E}_{x \sim p_{\text{data}}}[\log D(x)] + \mathbb{E}_{z \sim p_z}[\log(1 - D(G(z)))]$$
\end{formulabox}

\subsection{Effect of Regularization}
\begin{itemize}
    \item Prevents overfitting, encourages smooth/manifold latent spaces.
    \item In VAEs: KL divergence regularizes latent distribution toward standard normal, ensuring continuity (interpolation works).
    \item In GANs: Techniques like spectral normalization, gradient penalty stabilize training and improve sample diversity.
\end{itemize}

\subsection{Diffusion Models}
\begin{itemize}
    \item \textbf{Forward process}: Gradually add Gaussian noise to data over $T$ steps until it becomes pure noise.
    \item \textbf{Reverse process}: Train neural network to predict/remove noise at each step, recovering data from noise.
    \item \textbf{Pros}: Stable training, high diversity/fidelity, modular. \textbf{Cons}: Slow sampling (many steps required).
\end{itemize}

\subsection{Controlling Diffusion}
\begin{itemize}
    \item \textbf{Classifier Guidance}: Use gradients from a trained classifier to steer generation toward a specific class.
    \item \textbf{Classifier-Free Guidance}: Condition generation on text/labels during training. At inference, blend conditional and unconditional predictions. Higher $w$ = stronger conditioning.
\end{itemize}

\begin{formulabox}{Classifier-Free Guidance}
$$\hat{\epsilon} = \epsilon_u + w \cdot (\epsilon_c - \epsilon_u)$$
\begin{itemize}
    \item $\epsilon_c$ = conditional noise prediction, $\epsilon_u$ = unconditional
    \item $w > 1$ strengthens conditioning (e.g., $w = 7.5$ typical)
    \item $w = 1$ $\rightarrow$ purely conditional, $w = 0$ $\rightarrow$ unconditional
\end{itemize}
\end{formulabox}

\subsection{Flow Models (Flow Matching)}
\begin{itemize}
    \item Learns a continuous velocity field (ODE) that transports noise distribution to data distribution.
    \item \textbf{Advantage over Diffusion}: Direct trajectory, fewer steps needed, more efficient sampling.
\end{itemize}

\begin{formulabox}{Flow Matching ODE}
$$\frac{d\mathbf{x}}{dt} = v_\theta(\mathbf{x}, t)$$
Solve from noise $\mathbf{x}(0) \sim \mathcal{N}(0, 1)$ to data $\mathbf{x}(1)$ using numerical integrator (e.g., Euler method: $\mathbf{x}_{t+\Delta t} = \mathbf{x}_t + v_\theta(\mathbf{x}_t, t) \cdot \Delta t$).
\end{formulabox}

\section{Quick Reference Summary}
\begin{formulabox}{Must-Memorize Transform Matrices}
\begin{tabular}{>{\bfseries}p{2.5cm} p{10cm}}
Scale & $\begin{bmatrix} s_x & 0 & 0 \\ 0 & s_y & 0 \\ 0 & 0 & 1 \end{bmatrix}$ \\[1em]
Rotate ($\theta$) & $\begin{bmatrix} \cos\theta & -\sin\theta & 0 \\ \sin\theta & \cos\theta & 0 \\ 0 & 0 & 1 \end{bmatrix}$ \\[1em]
Translate & $\begin{bmatrix} 1 & 0 & t_x \\ 0 & 1 & t_y \\ 0 & 0 & 1 \end{bmatrix}$ \\[1em]
Shear & $\begin{bmatrix} 1 & a & 0 \\ b & 1 & 0 \\ 0 & 0 & 1 \end{bmatrix}$ \\[1em]
Reflect $y$-axis & $\begin{bmatrix} -1 & 0 & 0 \\ 0 & 1 & 0 \\ 0 & 0 & 1 \end{bmatrix}$
\end{tabular}
\end{formulabox}

\begin{formulabox}{Key Equations at a Glance}
\begin{tabular}{>{\bfseries}p{2.5cm} p{10cm}}
Blinn-Phong & $I = k_a I_a + k_d I_l (\vec{N}\!\cdot\!\vec{L}) + k_s I_l (\vec{N}\!\cdot\!\vec{H})^n$ \\[0.6em]
Halfway vector & $\vec{H} = \frac{\vec{L} + \vec{V}}{\|\vec{L} + \vec{V}\|}$ \\[0.8em]
NeRF color & $C = \sum T_i \alpha_i \mathbf{c}_i, \;\; T_i = \prod(1-\alpha_j)$ \\[0.8em]
Opacity & $\alpha_i = 1 - \exp(-\sigma_i \delta_i)$ \\[0.6em]
Positional encoding & $[\sin(2^k \pi \mathbf{x}), \cos(2^k \pi \mathbf{x})]$ \\[0.6em]
CFG & $\hat{\epsilon} = \epsilon_u + w(\epsilon_c - \epsilon_u)$ \\[0.6em]
Ray & $\mathbf{r}(t) = \mathbf{o} + t \cdot \mathbf{d}$ \\[0.6em]
Bézier & $P(t) = \sum B_{i,n}(t) P_i$ \\[0.6em]
Epipolar & $\mathbf{x}'^T F \mathbf{x} = 0$
\end{tabular}
\end{formulabox}

\end{document}